% Generated by tools/build_prior_work.py from the claims ledger; do not edit.
\bibitem{pwanderson2003introduction}
Anderson, T. W. (2003). \emph{An Introduction to Multivariate Statistical Analysis}. Wiley. ISBN 0-471-36091-0.

\bibitem{pwgraham1994concrete}
Graham, R. L., Knuth, D. E., and Patashnik, O. (1994). \emph{Concrete Mathematics: A Foundation for Computer Science}. Addison-Wesley. ISBN 0-201-55802-5.

\bibitem{pwhardy2008introduction}
Hardy, G. H., and Wright, E. M. (2008). \emph{An Introduction to the Theory of Numbers}. Oxford University Press. ISBN 978-0-19-921986-5.

\bibitem{pwhawking1973large}
Hawking, S. W., and Ellis, G. F. R. (1973). \emph{The Large Scale Structure of Space-Time}. Cambridge University Press. \url{https://doi.org/10.1017/CBO9780511524646}

\bibitem{pwhawking1976new}
Hawking, S. W., King, A. R., and McCarthy, P. J. (1976). ``A new topology for curved space-time which incorporates the causal, differential, and conformal structures.'' \emph{Journal of Mathematical Physics} \textbf{17}, 174--181. \url{https://doi.org/10.1063/1.522874}

\bibitem{pwhorn2012matrix}
Horn, R. A., and Johnson, C. R. (2012). \emph{Matrix Analysis}. Cambridge University Press. \url{https://doi.org/10.1017/CBO9781139020411}

\bibitem{pwmalament1977class}
Malament, D. B. (1977). ``The class of continuous timelike curves determines the topology of spacetime.'' \emph{Journal of Mathematical Physics} \textbf{18}, 1399--1404. \url{https://doi.org/10.1063/1.523436}

\bibitem{pwstein1971introduction}
Stein, E. M., and Weiss, G. (1971). \emph{Introduction to Fourier Analysis on Euclidean Spaces}. Princeton University Press (Princeton Mathematical Series). \url{https://doi.org/10.1515/9781400883899}

\bibitem{pwszego1939orthogonal}
Szeg{\H{o}}, G. (1939). \emph{Orthogonal Polynomials}. American Mathematical Society (Colloquium Publications 23). \url{https://doi.org/10.1090/coll/023}

